\documentclass[../../../include/open-logic-section]{subfiles}

\begin{document}

\olfileid{sth}{ord-arithmetic}{using-addition}
\olsection{ক্রমসংখ্যার যোগের প্রয়োগ}

ক্রমসংখ্যার যোগ ব্যবহার করে বিভিন্ন সেটের স্তরাঙ্ক
সরাসরি হিসাব করা যায়;
এখানে স্তরাঙ্ক \olref[spine][rank]{defnsetrank}-এর অর্থে:

\begin{lem}\ollabel{rankcomputation}
$\setrank{A} = \alpha$ এবং $\setrank{B} = \beta$ হলে:
\begin{enumerate}
	\item\ollabel{exrankpow} $\setrank{\Pow{A}} = \alpha\ordplus 1$;
	\item\ollabel{exrankpair} $\setrank{\{A, B\}} = \max(\alpha,
	\beta) \ordplus 1$;
	\item\ollabel{exrankcup} $\setrank{A \cup B} = \max(\alpha,
	\beta)$;
	\item\ollabel{exranktuple} $\setrank{\tuple{A,B}} = \max(\alpha,
	\beta) \ordplus 2$;
	\item\ollabel{exranktimes} $\setrank{A \times B} \leq \max(\alpha,
	\beta) \ordplus 2$;
	\item\ollabel{exrankunion} $\alpha$ শূন্য অথবা সীমা ক্রমসংখ্যা হলে
	$\setrank{\bigcup A} = \alpha$; $\alpha = \gamma\ordplus 1$
	হলে $\setrank{\bigcup A} = \gamma$।
\end{enumerate}
\end{lem}

\begin{proof}
সব ক্ষেত্রেই \olref[spine][rank]{ranksupstrict} বারবার
ব্যবহার করব।

\emph{\olref{exrankpow}।} $x \subseteq A$ হলে
$\setrank{x} \leq \setrank{A}$। তাই
$\setrank{\Pow{A}} \leq \alpha \ordplus 1$। বিশেষত
$A \in \Pow{A}$ হওয়ায় $\setrank{\Pow{A}} =
\alpha \ordplus 1$।

\emph{\olref{exrankpair}।} \olref[spine][rank]{ranksupstrict}
থেকে।

\emph{\olref{exrankcup}।} \olref[spine][rank]{ranksupstrict}
থেকে।

\emph{\olref{exranktuple}।} \olref{exrankpair} দুবার
প্রয়োগ করে।

% BN-SRC-449: the displayed double power-set bound uses exrankpow twice, not exranktuple.
\emph{\olref{exranktimes}।} লক্ষ করুন,
$A \times B \subseteq \Pow{\Pow{A \cup B}}$;
এবার \olref{exrankpow} দুবার প্রয়োগ করুন।

\emph{\olref{exrankunion}।} $\alpha = \gamma\ordplus 1$
হলে কোনো $c \in A$ আছে, যার
$\setrank{c} = \gamma$; আর $A$-র কোনো !!{element}-এর
স্তরাঙ্ক এর চেয়ে বেশি নয়। তাই
$\setrank{\bigcup A} = \gamma$। $\alpha$ সীমা
ক্রমসংখ্যা হলে $A$-তে এমন !!{element}s আছে, যাদের
স্তরাঙ্ক $\alpha$-র যথেচ্ছ কাছে যায় কিন্তু তার চেয়ে
কম থাকে। ফলে $\bigcup A$-তেও একইভাবে
!!{element}s-এর স্তরাঙ্ক $\alpha$-র যথেচ্ছ কাছে
যায় কিন্তু তার চেয়ে কম থাকে; তাই
$\setrank{\bigcup A} = \alpha$।
\end{proof}
\noindent
\olref{exranktimes}-এ সমতার বদলে কেন \emph{অসমতা}
আছে, তা দেখার কাজটি অনুশীলনের জন্য রইল।

\begin{prob}
% BN-SRC-450: the source exercise omits the equality sign before the upper bound.
এমন সেট $A$ ও $B$ দিন, যাতে $\setrank{A \times B} =
\max(\setrank{A}, \setrank{B})$। আবার এমন সেট
$A$ ও $B$ দিন, যাতে
$\setrank{A \times B}=\max(\setrank{A}, \setrank{B})
\ordplus 2$। অন্য কোনো স্তরাঙ্ক কি সম্ভব?
\end{prob}

এখন ``সসীম'' ও ``অসীম'' ক্রমসংখ্যা বলতে কী বোঝায়,
তার কয়েকটি স্বাভাবিক সংজ্ঞার সমতুল্যতাও দেখাতে পারি:

\begin{lem}\ollabel{ordinfinitycharacter}
যেকোনো ক্রমসংখ্যা $\alpha$-র জন্য নিচের দাবিগুলি সমতুল্য:
\begin{enumerate}
	\item\ollabel{ord:notinomega} $\alpha\notin \omega$,
	অর্থাৎ $\alpha$ স্বাভাবিক সংখ্যা নয়;
	\item\ollabel{ord:omegaplus} $\omega \leq \alpha$;
	\item\ollabel{ord:oneplus} $1 \ordplus \alpha = \alpha$;
	%\alpha \approx \alpha \ordplus 1$
	\item\ollabel{ord:plusone} $\alpha \approx \alpha\ordplus 1$,
	অর্থাৎ $\alpha$ ও $\alpha\ordplus 1$-এর উপাদানসংখ্যা সমান;
	\item\ollabel{ord:infinite} $\alpha$ দেদেকিন্ড-অসীম।
\end{enumerate}
\end{lem}
\noindent
অতএব একটি ক্রমসংখ্যার (অ)সসীম হওয়ার পাঁচটি
প্রমাণযোগ্যভাবে সমতুল্য ব্যাখ্যা পেলাম।

\begin{proof}
\emph{\olref{ord:notinomega} $\Rightarrow$
\olref{ord:omegaplus}।} ত্রিবিকল্প নীতি থেকে।

\emph{\olref{ord:omegaplus} $\Rightarrow$
\olref{ord:oneplus}।} $\alpha \geq \omega$ ধরি।
ট্রান্সফাইনাইট আবেশে দেখা যায়, এমন একটি ক্ষুদ্রতম
ক্রমসংখ্যা $\gamma$ (সম্ভবত $0$) আছে, যাতে কোনো
সীমা ক্রমসংখ্যা $\beta$-র জন্য
$\alpha = \beta \ordplus \gamma$। এবার:
\[
	1 \ordplus \alpha =
	1 \ordplus (\beta \ordplus \gamma) =
	(1 \ordplus \beta) \ordplus \gamma =
	\supstrict_{\delta < \beta} (1 \ordplus \delta) \ordplus \gamma =
	\beta \ordplus \gamma =
	\alpha.
\]
\emph{\olref{ord:oneplus} $\Rightarrow$
\olref{ord:plusone}।} স্পষ্টতই
$f \colon (\alpha \disjointsum 1) \to
(1 \disjointsum \alpha)$ একটি !!a{bijection}।
$1 \ordplus \alpha = \alpha$ হলে
$g \colon (1 \disjointsum \alpha) \to \alpha$
একটি সমরূপতা আছে। এবার $\comp{f}{g}$ বিবেচনা করুন।

\emph{\olref{ord:plusone} $\Rightarrow$
\olref{ord:infinite}।} $\alpha \approx \alpha \ordplus 1$
হলে $f \colon (\alpha \disjointsum 1) \to \alpha$
একটি !!a{bijection} আছে। প্রত্যেক $\gamma < \alpha$-র
জন্য $g(\gamma) = f(\gamma, 0)$ নিই। এই
!!{injection} দেখায় $\alpha$ দেদেকিন্ড-অসীম,
কারণ $f(0,1) \in \alpha \setminus \ran{g}$।

\emph{\olref{ord:infinite} $\Rightarrow$
\olref{ord:notinomega}।} এটি
\olref[z][infinity-again]{naturalnumbersarentinfinite}।
\end{proof}

\end{document}
